%-*-latex-*-

\documentclass[a4paper,11pt]{article}

% Language and fonts
%
\usepackage[utf8]{inputenc}
\usepackage[T1]{fontenc}
\usepackage[spanish]{babel}
\usepackage{hyphenat}
\usepackage{url}

\pagestyle{empty}

\begin{document}

\noindent
\begin{tabular}{@{}l@{\qquad\qquad\qquad}r@{}}
  Christian \textsc{Rinderknecht}
  & \\
  Széchenyi István utca, 8
  & \\
  3300 Eger, Húngria \\
  \texttt{+36 70 311 6130} \\
  \url{rinderknecht@free.fr}
\end{tabular}

\bigskip

\begin{center}
  \large Carta de presentación
\end{center}

\bigskip

Después de mi doctorado en informática, trabajé como ingeniero en
empresas, pero también como profesor en varias universidades hasta
\oldstylenums{2014} (ELTE en Budapest fue la última). Durante esos
años, enseñé en inglés, en varios países; aprendí a conocer muchas
culturas; publiqué miles de páginas en inglés, y cientos en francés.


En paralelo a mi carrera en informática, cultivé mi pasión por la
traducción de poesía, y empecé con Pablo Neruda en \oldstylenums{1998}
(editorial Gallimard). Desde entonces publiqué varias traducciones, en
mayoría del francés al español y viceversa, pero también, este año,
del húngaro hacia el francés, el inglés y el español (poemas de Attila
József). Aún no hablo el húngaro, y lo voy a estudiar más. Por otro
lado, también publiqué mis propios poemas en inglés.

Deseo cambiar de profesión y volver a enseñar, esta vez idiomas y
literatura. Enseñé el francés y el español en privado, y me gustaría
hoy día hacer parte de una escuela de idiomas.

Tengo una empresa en Húngria (Kft) y puedo facturar mis servicios, tan
como podría convertirme en su empleado.

\bigskip
\bigskip

\hfill Hecho en Eger, el \oldstylenums{28} Abril \oldstylenums{2026}.

\bigskip
\bigskip
\bigskip
\bigskip
\hfill Christian Rinderknecht

\bigskip

\end{document}
